% GENERATED FILE. Edit canonical lessons, then run npm run generate:books.
% canonical-source-hash: fnv1a64:9538615432185f0f
% canonical-lessons: GE-C479-erkaelten, GE-C479-schnupfen, GE-C479-hustensaft, GE-C479-schweiss, GE-C479-gaensehaut

\chapter{Schnupfen und Gänsehaut}
\label{ch:ge-schnupfen-und-gansehaut}
\hlchaptermodality{479}

\begin{chapteropening}
\textbf{By the end of this chapter:} \emph{I can talk about colds, sweat and shivers.}

Complete the last lesson of chapter 479: I can talk about colds, sweat and shivers.
\end{chapteropening}

\section[erkälten]{erkälten — to catch a cold (sich erkälten)}

\label{lesson:GE-C479-erkaelten}



\begin{quote}
Before the new one: say the German for to insult, then the German for to hurt.
\end{quote}

\subsection*{You'll want to know: erkälten}
\textbf{erkälten} — \textquotedblleft{}to catch a cold (sich erkälten)\textquotedblright{}.

\begin{quote}
Im Regen habe ich mich erkältet.
\end{quote}

\begin{quote}
\emph{I caught a cold in the rain.}
\end{quote}

Reflexive: \emph{sich erkälten}. \emph{Erkältet sein} — to have a cold.

\begin{cousinweb}
\emph{erkälten} comes from Middle High German \emph{erkelten} (EtymDB), built on \emph{kalt} \textquotedblleft{}cold\textquotedblright{}, which is inherited from Old High German \emph{kalt}, from Proto-Germanic \emph{kaldaz}.
\end{cousinweb}

\begin{grammarlens}[title={What you've built}]
One more word for telling, explaining and describing at B1.
\end{grammarlens}

\subsection*{Your turn}
\begin{itemize}
\raggedright
  \item \emph{Say it:} \emph{erkälten}
  \item \emph{Say it:} \emph{erkälten}, once more
  \item \emph{Say it:} a sentence of your own with \emph{erkälten}
  \item \emph{Recall:} say the German for to insult, then the German for to hurt, then say \emph{erkälten} again
\end{itemize}

\begin{tcolorbox}[breakable,colback=teal!4,colframe=teal!35!black,title={Before you move on}]
What does erkälten mean? (\textquotedblleft{}To catch a cold (sich erkälten)\textquotedblright{}.) Say it once more.
\end{tcolorbox}
\section[der Schnupfen]{der Schnupfen — head cold, runny nose}

\label{lesson:GE-C479-schnupfen}



\begin{quote}
Before the new one: say the German for to hurt, then the German for to catch a cold.
\end{quote}

\subsection*{You'll want to know: der Schnupfen}
\textbf{der Schnupfen} — \textquotedblleft{}head cold, runny nose\textquotedblright{}.

\begin{quote}
Ich habe Schnupfen und brauche viele Taschentücher.
\end{quote}

\begin{quote}
\emph{I have a runny nose and need lots of tissues.}
\end{quote}

Usually without an article: \emph{Ich habe Schnupfen.} No plural in everyday use.

\begin{cousinweb}
\emph{Schnupfen} is built on the verb \emph{schnupfen} \textquotedblleft{}to sniff\textquotedblright{} (EtymDB), probably an imitative word for the sound of a blocked nose.
\end{cousinweb}

\begin{grammarlens}[title={What you've built}]
One more word for telling, explaining and describing at B1.
\end{grammarlens}

\subsection*{Your turn}
\begin{itemize}
\raggedright
  \item \emph{Say it:} \emph{der Schnupfen}
  \item \emph{Say it:} \emph{der Schnupfen}, once more
  \item \emph{Say it:} a sentence of your own with \emph{der Schnupfen}
  \item \emph{Recall:} say the German for to hurt, then the German for to catch a cold, then say \emph{der Schnupfen} again
\end{itemize}

\begin{tcolorbox}[breakable,colback=teal!4,colframe=teal!35!black,title={Before you move on}]
What does der Schnupfen mean? (\textquotedblleft{}Head cold, runny nose\textquotedblright{}.) Say it once more.
\end{tcolorbox}
\section[der Hustensaft]{der Hustensaft — cough syrup}

\label{lesson:GE-C479-hustensaft}



\begin{quote}
Before the new one: say the German for to catch a cold, then the German for head cold.
\end{quote}

\subsection*{You'll want to know: der Hustensaft}
\textbf{der Hustensaft} — \textquotedblleft{}cough syrup\textquotedblright{}.

\begin{quote}
Der Hustensaft schmeckt nach Kirschen.
\end{quote}

\begin{quote}
\emph{The cough syrup tastes of cherries.}
\end{quote}

Plural: \emph{die Hustensäfte}. Built on \emph{der Husten} — cough, and \emph{der Saft} — juice.

\begin{cousinweb}
\emph{Hustensaft} is built on \emph{Husten} \textquotedblleft{}cough\textquotedblright{}, from Old High German \emph{huosto}, and \emph{Saft} \textquotedblleft{}juice\textquotedblright{}, from Middle High German \emph{saf} (EtymDB).
\end{cousinweb}

\begin{grammarlens}[title={What you've built}]
One more word for telling, explaining and describing at B1.
\end{grammarlens}

\subsection*{Your turn}
\begin{itemize}
\raggedright
  \item \emph{Say it:} \emph{der Hustensaft}
  \item \emph{Say it:} \emph{der Hustensaft}, once more
  \item \emph{Say it:} a sentence of your own with \emph{der Hustensaft}
  \item \emph{Recall:} say the German for to catch a cold, then the German for head cold, then say \emph{der Hustensaft} again
\end{itemize}

\begin{tcolorbox}[breakable,colback=teal!4,colframe=teal!35!black,title={Before you move on}]
What does der Hustensaft mean? (\textquotedblleft{}Cough syrup\textquotedblright{}.) Say it once more.
\end{tcolorbox}
\section[der Schweiß]{der Schweiß — sweat}

\label{lesson:GE-C479-schweiss}



\begin{quote}
Before the new one: say the German for head cold, then the German for cough syrup.
\end{quote}

\subsection*{You'll want to know: der Schweiß}
\textbf{der Schweiß} — \textquotedblleft{}sweat\textquotedblright{}.

\begin{quote}
Nach dem Joggen war mein T-Shirt nass vor Schweiß.
\end{quote}

\begin{quote}
\emph{After jogging my T-shirt was wet with sweat.}
\end{quote}

No plural. The verb is \emph{schwitzen} — to sweat.

\begin{cousinweb}
\emph{Schweiß} is inherited from Middle High German \emph{sweiz} and Old High German \emph{sweiz}, from Proto-Germanic \emph{swait-} \textquotedblleft{}sweat\textquotedblright{} (EtymDB); English \emph{sweat} is a cousin.
\end{cousinweb}

\begin{grammarlens}[title={What you've built}]
One more word for telling, explaining and describing at B1.
\end{grammarlens}

\subsection*{Your turn}
\begin{itemize}
\raggedright
  \item \emph{Say it:} \emph{der Schweiß}
  \item \emph{Say it:} \emph{der Schweiß}, once more
  \item \emph{Say it:} a sentence of your own with \emph{der Schweiß}
  \item \emph{Recall:} say the German for head cold, then the German for cough syrup, then say \emph{der Schweiß} again
\end{itemize}

\begin{tcolorbox}[breakable,colback=teal!4,colframe=teal!35!black,title={Before you move on}]
What does der Schweiß mean? (\textquotedblleft{}Sweat\textquotedblright{}.) Say it once more.
\end{tcolorbox}
\section[die Gänsehaut]{die Gänsehaut — goose bumps}

\label{lesson:GE-C479-gaensehaut}



\begin{quote}
Before the new one: say the German for cough syrup, then the German for sweat.
\end{quote}

\subsection*{You'll want to know: die Gänsehaut}
\textbf{die Gänsehaut} — \textquotedblleft{}goose bumps\textquotedblright{}.

\begin{quote}
Bei dieser Musik bekomme ich Gänsehaut.
\end{quote}

\begin{quote}
\emph{This music gives me goose bumps.}
\end{quote}

No plural. As in English, built on the goose: \emph{die Gans}.

\begin{cousinweb}
\emph{Gänsehaut} is built on \emph{Gans} \textquotedblleft{}goose\textquotedblright{}, inherited from Old High German \emph{gans}, from Proto-Germanic \emph{gans} (EtymDB), and \emph{Haut} \textquotedblleft{}skin\textquotedblright{}: skin like a plucked goose.
\end{cousinweb}

\begin{grammarlens}[title={What you've built}]
One more word for telling, explaining and describing at B1.
\end{grammarlens}

\subsection*{Your turn}
\begin{itemize}
\raggedright
  \item \emph{Say it:} \emph{die Gänsehaut}
  \item \emph{Say it:} \emph{die Gänsehaut}, once more
  \item \emph{Say it:} a sentence of your own with \emph{die Gänsehaut}
  \item \emph{Recall:} say the German for cough syrup, then the German for sweat, then say \emph{die Gänsehaut} again
\end{itemize}

\begin{tcolorbox}[breakable,colback=teal!4,colframe=teal!35!black,title={Before you move on}]
What does die Gänsehaut mean? (\textquotedblleft{}Goose bumps\textquotedblright{}.) Say it once more.
\end{tcolorbox}
